\chapter{Herpes Zoster (Shingles)}

\section*{Definition and Overview}
Herpes zoster, commonly known as shingles, is a painful blistering rash caused by reactivation of the varicella-zoster virus, the same virus that causes chickenpox. After a chickenpox infection, the virus lies dormant in nerve roots and can reactivate years later, usually when immunity declines with age or illness. Shingles typically affects one side of the body with a band of painful blisters. The most feared complication is post-herpetic neuralgia, persistent nerve pain after the rash heals. Vaccination is highly effective at preventing shingles and its complications.

\section*{Epidemiology}
Around one in three people who have had chickenpox will develop shingles in their lifetime, which equates to roughly 50,000 cases in Australia each year. The risk rises steeply with age: most cases occur in people over 50, and it is uncommon in children and young adults. Approximately 10--18\% of people with shingles develop post-herpetic neuralgia, with the risk increasing with age. The shingles vaccine is recommended for all adults over 50 and for immunocompromised people over 18, and vaccination rates are increasing.

\section*{Aetiology and Pathophysiology}
\begin{itemize}
    \item \textbf{Cause:} reactivation of varicella-zoster virus, which persists in the dorsal root ganglia after chickenpox
    \item \textbf{Trigger:} declining cell-mediated immunity with age, illness, stress, or immunosuppressive treatment
    \item \textbf{Pathophysiology:} the reactivated virus travels down sensory nerves to the skin, causing inflammation of the nerve and the characteristic rash in the territory of that nerve
    \item \textbf{Distribution:} typically a single dermatome on one side of the trunk or face
    \item \textbf{Infectivity:} the rash is contagious to people who have not had chickenpox or vaccination, causing chickenpox, not shingles
    \item \textbf{Post-herpetic neuralgia:} nerve damage from inflammation causes persistent pain after healing
\end{itemize}

\section*{Clinical Features}
\begin{itemize}
    \item Burning, tingling, or shooting pain in a band on one side of the body, often before any rash
    \item A red rash that develops into clusters of fluid-filled blisters
    \item The rash follows a nerve distribution, usually on the chest, abdomen, or face
    \item Pain, itching, and sensitivity in the affected area
    \item Blisters crust over and heal within 2--4 weeks
    \item Fever, headache, and fatigue may accompany the rash
    \item HZO: shingles involving the eye (herpes zoster ophthalmicus) causes eye pain and threatens vision
\end{itemize}

\section*{Differential Diagnosis}
\begin{itemize}
    \item Herpes simplex infection, which has a similar blistering appearance
    \item Contact dermatitis and allergic reactions
    \item Cellulitis and other bacterial skin infections
    \item Chest pain from cardiac or lung causes, before the rash appears
    \item Trigeminal neuralgia and other facial pain
    \item The unilateral, dermatomal distribution and blistering are characteristic
\end{itemize}

\section*{When to Seek Medical Care}
Anyone who suspects shingles should see a doctor promptly, ideally within 72 hours of the rash appearing, because antiviral treatment is most effective when started early. Shingles on the face, especially near the eye, requires urgent review.

\textbf{Call 000 (or local emergency services) for:}
\begin{itemize}
    \item Shingles involving the eye, with pain, redness, or vision changes
    \item Rash with severe headache, confusion, or neurological symptoms
    \item Shingles in someone with a severely weakened immune system
    \item Widespread rash with fever and feeling very unwell
    \item Signs of secondary bacterial infection: spreading redness, warmth, and pus
\end{itemize}

\section*{Diagnosis and Investigations}
\begin{itemize}
    \item \textbf{Clinical diagnosis:} the characteristic unilateral rash in a nerve distribution is usually diagnostic
    \item \textbf{PCR swab:} confirms the diagnosis in atypical cases
    \item \textbf{Blood tests:} rarely needed, but antibody levels can support the diagnosis
    \item \textbf{Eye assessment:} urgent slit-lamp examination for suspected eye involvement
    \item \textbf{Assessment of immune status:} relevant in recurrent or severe disease
\end{itemize}

\section*{Management}
\begin{itemize}
    \item \textbf{Antiviral treatment:} aciclovir, valaciclovir, or famciclovir, started within 72 hours, speeds healing and reduces post-herpetic neuralgia
    \item \textbf{Pain relief:} simple analgesics, and stronger options for severe pain
    \item \textbf{Skin care:} keep the rash clean and dry; cool compresses soothe itching
    \item \textbf{Post-herpetic neuralgia treatment:} specialist pain medicines, nerve blocks, and capsaicin or lidocaine patches
    \item \textbf{Eye involvement:} urgent ophthalmology care with antiviral and topical treatment
    \item \textbf{Prevention of spread:} cover the rash and avoid contact with people who have not had chickenpox
    \item \textbf{Vaccination:} shingles vaccine is given after recovery to prevent future episodes
\end{itemize}

\section*{Complications}
\begin{itemize}
    \item Post-herpetic neuralgia, persistent pain lasting months or years, the most common and disabling complication
    \item Herpes zoster ophthalmicus with corneal scarring and vision loss
    \item Secondary bacterial infection of the rash
    \item Ramsay Hunt syndrome: shingles affecting the ear and facial nerve, causing facial weakness and hearing problems
    \item Neurological complications, including encephalitis and myelitis, rarely
    \item Spread to people who have not had chickenpox
\end{itemize}

\section*{Prognosis and Outcomes}
Most people with shingles recover fully within a few weeks, with the rash healing and pain settling. Early antiviral treatment reduces both the severity of the illness and the risk of post-herpetic neuralgia. Post-herpetic neuralgia is more common in older people and can be challenging to treat, but many cases resolve within a year and modern pain management helps. Vaccination dramatically reduces the risk of shingles and its complications, making it one of the most effective preventive measures available for older adults.

\section*{Prevention and Self-Care}
\begin{itemize}
    \item Have the shingles vaccine, recommended for all adults over 50 and immunocompromised adults over 18
    \item Two doses of the recombinant shingles vaccine provide over 90\% protection
    \item Start antiviral treatment within 72 hours of the rash if shingles occurs
    \item Keep the rash covered to prevent spread to others
    \item Rest, eat well, and manage pain with regular analgesia
    \item Seek urgent care for facial or eye involvement
    \item Manage stress and stay well to reduce reactivation risk
    \item Discuss vaccination with your doctor, even after a previous episode of shingles
\end{itemize}

\section*{Sources and Further Reading}
The following reputable sources provide further information for healthcare students and patients seeking reliable, current advice about shingles.

\begin{itemize}
    \item Healthdirect. \textit{Shingles (herpes zoster).} https://www.healthdirect.gov.au/shingles
    \item Australian Government Department of Health. \textit{Shingles vaccination program.} https://www.health.gov.au/
    \item Shingles Australia and Pain Australia. \textit{Post-herpetic neuralgia information.} https://www.painaustralia.org.au/
    \item NHS. \textit{Shingles: overview and treatment.} https://www.nhs.uk/conditions/shingles/
    \item Mayo Clinic. \textit{Shingles: symptoms and causes.} https://www.mayoclinic.org/
    \item Centers for Disease Control and Prevention. \textit{Shingles (herpes zoster).} https://www.cdc.gov/
\end{itemize}
